\scanpage{151}\mspage{25}
\[
\left\{\begin{aligned}
&\text{les }M^\pi=\underline{\mathrm{Filt}}^\pi\text{ sont stables par sommes amalgamées}\\
&\text{simples, et par }\underset{\unsure{\mathrel{\text{---}}|\longrightarrow}}{\lim}\text{ filtrantes de cardinal}\\
&\le\pi\quad\text{(pas néc. \unsure{grandes|grand} dev. \unsure{à|de} }\pi_0),
\end{aligned}\right.
\tag*{(3.13)}
\]
compte tenu que $r\le s$ donc $M^r\subset M^s$, donc $A,B,X$ \unsure{donc|dans} $X_D=B\amalg_A X$ est dans $M^s$. Les sommes amalgamées partielles \uline{finies} des $X_D$ sont donc aussi dans $M^s$, par récurrence sur le cardinal de l'ens d'indices\unsure{;|,} \deleted{\unsure{étant|et d'}} La somme amalgamée infinie apparaît ainsi comme $\underset{\longrightarrow}{\lim}$ filtrante d'objets de $M^s$, l'ens d'indices étant l'ens. $\unsure{\mathscr{P}|\mathcal{P}}^*_f\Delta\unsure{(}C_0,f\unsure{)}$ des parties finies non vides de $\Delta(C_0,f)$. Cet ens. est de cardinal majoré par \deleted{$\mathrm{card}\,\Delta(C_0,f)\cdot\unsure{s|2|8}$} $s'$, \deleted{\unsure{dès|ds}} \deleted{\unsure{pour|que}} si $s'$ est un cardinal infini tel que $\mathrm{card}\,\Delta(C_0,f)\le s'$. On peut donc prendre $s'=s^{\unsure{(}r^{\pi_0})}$.
On a donc prouvé, compte tenu de (3.13) appliqué à $M^{s'}$, le

\uline{Lemme} 3.3. On suppose (3.10) i.e. $C_0\subset\mathrm{Fl}^rM$, $f\in\mathrm{Fl}^sM$, avec \deleted{\unsure{assez}} \inserted{$s\ge r\ge\mathrm{Sup}(\pi_1,C_1)$,} et de plus que $C_0$ satisfait la condition de non pléthoricité (3.11.) (ce qu'on peut toujours supposer, quitte à réduire $C_0$ sans changer $\widetilde{C}_0=\Phi$ et $C_{0\unsure{\alpha|*}}=\Psi$), et que la filtration cardinale satisfait la \unsure{condition|conditions} de stabilité (\unsure{condition|commodité|commode} !) (3.13.). Posons
\[
s'=s^{\unsure{(}r^{\pi_0})}.
\]
Alors $X_1(C_0,f)$ est dans $M^{s'}$. Plus \unsure{généralement} et plus \unsure{précisément}, pour $C_0$ et $r$ fixés (avec $C_0\subset\mathrm{Fl}^rM$ et satisfaisant la condition de non pléthoricité (3.11) relativement à $r$), pour tout cardinal
% archive sheet number: 150
